\documentclass[10pt]{extarticle}
\usepackage[a4paper,margin=1.5cm]{geometry}
\usepackage[T1]{fontenc}
\usepackage{amsmath,amssymb}
\usepackage{graphicx}
\usepackage[export]{adjustbox}
\usepackage{multicol}
\usepackage{enumitem}
\usepackage{booktabs,array,tabularx}
\setlength{\columnsep}{0.7cm}
\setlength{\parindent}{0pt}
\setlength{\parskip}{3pt}
\pagestyle{empty}
\makeatletter
% figure registry (figdefs.tex is generated by the extract step)
\newcommand{\qfigdef}[4]{\@namedef{qf@file@#1}{#2}\@namedef{qf@w@#1}{#3}\@namedef{qf@wide@#1}{#4}}
\newcommand{\qf@check}[1]{\@ifundefined{qf@file@#1}{\PackageError{score2}{Unknown figure id #1}{}}{}}
\newcommand{\qf@img}[4][]{\includegraphics[width=#2\dimexpr\@nameuse{qf@w@#4}pt\relax,max width=#3,max height=0.45\textheight,#1]{\@nameuse{qf@file@#4}}}
% Full-width questions: finalize wraps every question that contains a wide figure in \qfullbegin...\qfullend,
% so the WHOLE question (stem, figures, options) is set across the page in reading order, then the two columns
% resume. The question's head - its number, stem and first figure (\qfullsplit marks the end of the head) -
% is kept on one page, so a number is never left at a page foot with its figure overleaf; the rest may break
% across pages normally, so no big gaps are left. Without a \qfullsplit, the whole question is kept together
% if it is under 45% of a page. Never write these yourself.
\newif\ifqf@full
\newif\ifqf@boxed
\newsavebox{\qf@box}
\newcommand{\qf@open}{\global\qf@fulltrue\par\global\qf@boxedtrue
  \setbox\qf@box\vbox\bgroup\hsize\textwidth\linewidth\textwidth}
\newcommand{\qf@emit}[1]{\par\egroup\global\qf@boxedfalse
  \ifdim\dimexpr\ht\qf@box+\dp\qf@box\relax<#1\textheight\noindent\usebox{\qf@box}\par\else\unvbox\qf@box\fi}
\newcommand{\qfullsplit}{\ifqf@boxed\qf@emit{0.85}\fi}
\newcommand{\qf@close}{\ifqf@boxed\qf@emit{0.45}\fi\global\qf@fullfalse}
\newcommand{\qfullbegin}{\end{multicols}\qf@open}
\newcommand{\qfullend}{\qf@close\begin{multicols}{2}}
\newcommand{\qfullstart}{\qf@open}
\newcommand{\qfulldone}{\qf@close}
% \qfig[scale]{id}: block figure, centred, directly where it is written (size set by the tool).
% Fallback if a wide figure is not inside a full-width question: break out of the columns in place - no floats.
\newcommand{\qfig}[2][1]{\qf@check{#2}%
  \ifqf@full\qf@blk{#1}{#2}\else
  \ifnum\@nameuse{qf@wide@#2}=1 \expandafter\@firstoftwo\else\expandafter\@secondoftwo\fi
  {\end{multicols}\begin{center}\qf@img{#1}{\textwidth}{#2}\end{center}\begin{multicols}{2}}%
  {\qf@blk{#1}{#2}}\fi}
\newcommand{\qf@blk}[2]{\par\smallskip{\centering\qf@img{#1}{\linewidth}{#2}\par}\smallskip}
% \qfigi[scale]{id}: inline figure (inside option lists / tables / minipages); never breaks columns.
\newcommand{\qfigi}[2][1]{\qf@check{#2}\qf@img[valign=c]{#1}{\linewidth}{#2}}
\makeatother
% structure macros - lint reads these to check completeness and figure placement
\newcommand{\papersec}[1]{\par\bigskip{\centering\Large\bfseries PAPER-#1\par}\medskip}
\newcommand{\partsec}[2]{\par\medskip\noindent{\large\bfseries PART-#1: #2}\par\smallskip}
\newcommand{\sect}[1]{\par\smallskip\noindent\textbf{#1}\par}
\newcommand{\qn}[1]{\par\smallskip\noindent\textbf{#1.}~}
\newcommand{\paraq}[2]{\par\smallskip\noindent\textbf{Paragraph for Questions #1 and #2}\par}
\newenvironment{qoptions}{\begin{itemize}[leftmargin=1.8em,itemsep=1pt,topsep=2pt,label={}]}{\end{itemize}}
\newcommand{\opt}[2]{\item[(#1)] #2}
% figure registry (generated by score2_tools.py finalize)
\qfigdef{p1_phy_s1_q01_1}{Nurture SRG - Test 03 (04 Feb 2024)-img/p1_phy_s1_q01_1.png}{72.2}{0}
\qfigdef{p1_phy_s1_q03_1}{Nurture SRG - Test 03 (04 Feb 2024)-img/p1_phy_s1_q03_1.png}{134.9}{0}
\qfigdef{p1_phy_s2_q08_1}{Nurture SRG - Test 03 (04 Feb 2024)-img/p1_phy_s2_q08_1.png}{127.1}{0}
\qfigdef{p1_phy_s2_q09_1}{Nurture SRG - Test 03 (04 Feb 2024)-img/p1_phy_s2_q09_1.png}{148.2}{0}
\qfigdef{p1_phy_s2_q12_1}{Nurture SRG - Test 03 (04 Feb 2024)-img/p1_phy_s2_q12_1.png}{152.5}{0}
\qfigdef{p1_phy_s3_q03_1}{Nurture SRG - Test 03 (04 Feb 2024)-img/p1_phy_s3_q03_1.png}{149.2}{0}
\qfigdef{p1_phy_s3_q06_1}{Nurture SRG - Test 03 (04 Feb 2024)-img/p1_phy_s3_q06_1.png}{114.6}{0}
\qfigdef{p1_chem_s1_q01_A}{Nurture SRG - Test 03 (04 Feb 2024)-img/p1_chem_s1_q01_A.png}{164.9}{0}
\qfigdef{p1_chem_s1_q01_B}{Nurture SRG - Test 03 (04 Feb 2024)-img/p1_chem_s1_q01_B.png}{187.6}{0}
\qfigdef{p1_chem_s1_q01_C}{Nurture SRG - Test 03 (04 Feb 2024)-img/p1_chem_s1_q01_C.png}{203.8}{0}
\qfigdef{p1_chem_s1_q01_D}{Nurture SRG - Test 03 (04 Feb 2024)-img/p1_chem_s1_q01_D.png}{35.1}{0}
\qfigdef{p1_chem_s1_q02_1}{Nurture SRG - Test 03 (04 Feb 2024)-img/p1_chem_s1_q02_1.png}{78.6}{0}
\qfigdef{p1_chem_s1_q02_2}{Nurture SRG - Test 03 (04 Feb 2024)-img/p1_chem_s1_q02_2.png}{78.2}{0}
\qfigdef{p1_chem_s1_q02_3}{Nurture SRG - Test 03 (04 Feb 2024)-img/p1_chem_s1_q02_3.png}{78.1}{0}
\qfigdef{p1_chem_s1_q02_4}{Nurture SRG - Test 03 (04 Feb 2024)-img/p1_chem_s1_q02_4.png}{79.8}{0}
\qfigdef{p1_chem_s1_q03_1}{Nurture SRG - Test 03 (04 Feb 2024)-img/p1_chem_s1_q03_1.png}{182.2}{0}
\qfigdef{p1_chem_s1_q03_A}{Nurture SRG - Test 03 (04 Feb 2024)-img/p1_chem_s1_q03_A.png}{114.2}{0}
\qfigdef{p1_chem_s1_q03_B}{Nurture SRG - Test 03 (04 Feb 2024)-img/p1_chem_s1_q03_B.png}{88.2}{0}
\qfigdef{p1_chem_s1_q03_C}{Nurture SRG - Test 03 (04 Feb 2024)-img/p1_chem_s1_q03_C.png}{105.6}{0}
\qfigdef{p1_chem_s1_q03_D}{Nurture SRG - Test 03 (04 Feb 2024)-img/p1_chem_s1_q03_D.png}{102.1}{0}
\qfigdef{p1_chem_s1_q04_1}{Nurture SRG - Test 03 (04 Feb 2024)-img/p1_chem_s1_q04_1.png}{102.6}{0}
\qfigdef{p1_chem_s1_q04_A}{Nurture SRG - Test 03 (04 Feb 2024)-img/p1_chem_s1_q04_A.png}{67.1}{0}
\qfigdef{p1_chem_s1_q04_B}{Nurture SRG - Test 03 (04 Feb 2024)-img/p1_chem_s1_q04_B.png}{102.7}{0}
\qfigdef{p1_chem_s1_q04_C}{Nurture SRG - Test 03 (04 Feb 2024)-img/p1_chem_s1_q04_C.png}{91.8}{0}
\qfigdef{p1_chem_s1_q04_D}{Nurture SRG - Test 03 (04 Feb 2024)-img/p1_chem_s1_q04_D.png}{103.5}{0}
\qfigdef{p1_chem_s1_q06_1}{Nurture SRG - Test 03 (04 Feb 2024)-img/p1_chem_s1_q06_1.png}{159.9}{0}
\qfigdef{p1_chem_s1_q06_A}{Nurture SRG - Test 03 (04 Feb 2024)-img/p1_chem_s1_q06_A.png}{80.1}{0}
\qfigdef{p1_chem_s1_q06_B}{Nurture SRG - Test 03 (04 Feb 2024)-img/p1_chem_s1_q06_B.png}{84.2}{0}
\qfigdef{p1_chem_s1_q06_C}{Nurture SRG - Test 03 (04 Feb 2024)-img/p1_chem_s1_q06_C.png}{80.1}{0}
\qfigdef{p1_chem_s1_q06_D}{Nurture SRG - Test 03 (04 Feb 2024)-img/p1_chem_s1_q06_D.png}{75.7}{0}
\qfigdef{p1_chem_s2_q07_A}{Nurture SRG - Test 03 (04 Feb 2024)-img/p1_chem_s2_q07_A.png}{96.2}{0}
\qfigdef{p1_chem_s2_q07_B}{Nurture SRG - Test 03 (04 Feb 2024)-img/p1_chem_s2_q07_B.png}{58.3}{0}
\qfigdef{p1_chem_s2_q07_C}{Nurture SRG - Test 03 (04 Feb 2024)-img/p1_chem_s2_q07_C.png}{75.1}{0}
\qfigdef{p1_chem_s2_q07_D}{Nurture SRG - Test 03 (04 Feb 2024)-img/p1_chem_s2_q07_D.png}{32.6}{0}
\qfigdef{p1_chem_s2_q08_1}{Nurture SRG - Test 03 (04 Feb 2024)-img/p1_chem_s2_q08_1.png}{132.7}{0}
\qfigdef{p1_chem_s2_q10_A}{Nurture SRG - Test 03 (04 Feb 2024)-img/p1_chem_s2_q10_A.png}{96.3}{0}
\qfigdef{p1_chem_s2_q10_B}{Nurture SRG - Test 03 (04 Feb 2024)-img/p1_chem_s2_q10_B.png}{101.1}{0}
\qfigdef{p1_chem_s2_q10_C}{Nurture SRG - Test 03 (04 Feb 2024)-img/p1_chem_s2_q10_C.png}{97.3}{0}
\qfigdef{p1_chem_s2_q10_D}{Nurture SRG - Test 03 (04 Feb 2024)-img/p1_chem_s2_q10_D.png}{78.9}{0}
\qfigdef{p1_chem_s2_q11_1}{Nurture SRG - Test 03 (04 Feb 2024)-img/p1_chem_s2_q11_1.png}{180.7}{0}
\qfigdef{p1_chem_s2_q11_2}{Nurture SRG - Test 03 (04 Feb 2024)-img/p1_chem_s2_q11_2.png}{78.7}{0}
\qfigdef{p1_chem_s2_q12_1}{Nurture SRG - Test 03 (04 Feb 2024)-img/p1_chem_s2_q12_1.png}{230.6}{0}
\qfigdef{p1_chem_s2_q12_2}{Nurture SRG - Test 03 (04 Feb 2024)-img/p1_chem_s2_q12_2.png}{96.7}{0}
\qfigdef{p1_chem_s2_q12_3}{Nurture SRG - Test 03 (04 Feb 2024)-img/p1_chem_s2_q12_3.png}{43.5}{0}
\qfigdef{p1_chem_s3_q01_1}{Nurture SRG - Test 03 (04 Feb 2024)-img/p1_chem_s3_q01_1.png}{205.4}{0}
\qfigdef{p1_chem_s3_q03_1}{Nurture SRG - Test 03 (04 Feb 2024)-img/p1_chem_s3_q03_1.png}{70.0}{0}
\qfigdef{p1_chem_s3_q03_2}{Nurture SRG - Test 03 (04 Feb 2024)-img/p1_chem_s3_q03_2.png}{48.7}{0}
\qfigdef{p1_chem_s3_q04_1}{Nurture SRG - Test 03 (04 Feb 2024)-img/p1_chem_s3_q04_1.png}{208.0}{0}
\qfigdef{p1_chem_s3_q05_1}{Nurture SRG - Test 03 (04 Feb 2024)-img/p1_chem_s3_q05_1.png}{207.5}{0}
\qfigdef{p1_chem_s3_q06_1}{Nurture SRG - Test 03 (04 Feb 2024)-img/p1_chem_s3_q06_1.png}{193.0}{0}
\qfigdef{p2_phy_s1_q02_1}{Nurture SRG - Test 03 (04 Feb 2024)-img/p2_phy_s1_q02_1.png}{72.5}{0}
\qfigdef{p2_phy_s1_q03_1}{Nurture SRG - Test 03 (04 Feb 2024)-img/p2_phy_s1_q03_1.png}{103.0}{0}
\qfigdef{p2_phy_s2_q05_1}{Nurture SRG - Test 03 (04 Feb 2024)-img/p2_phy_s2_q05_1.png}{89.6}{0}
\qfigdef{p2_phy_s2_q07_1}{Nurture SRG - Test 03 (04 Feb 2024)-img/p2_phy_s2_q07_1.png}{98.7}{0}
\qfigdef{p2_phy_s2_q12_1}{Nurture SRG - Test 03 (04 Feb 2024)-img/p2_phy_s2_q12_1.png}{149.2}{0}
\qfigdef{p2_phy_s3_q01_1}{Nurture SRG - Test 03 (04 Feb 2024)-img/p2_phy_s3_q01_1.png}{101.3}{0}
\qfigdef{p2_chem_s1_q01_A}{Nurture SRG - Test 03 (04 Feb 2024)-img/p2_chem_s1_q01_A.png}{155.1}{0}
\qfigdef{p2_chem_s1_q01_B}{Nurture SRG - Test 03 (04 Feb 2024)-img/p2_chem_s1_q01_B.png}{120.1}{0}
\qfigdef{p2_chem_s1_q01_C}{Nurture SRG - Test 03 (04 Feb 2024)-img/p2_chem_s1_q01_C.png}{137.0}{0}
\qfigdef{p2_chem_s1_q01_D}{Nurture SRG - Test 03 (04 Feb 2024)-img/p2_chem_s1_q01_D.png}{127.5}{0}
\qfigdef{p2_chem_s1_q02_1}{Nurture SRG - Test 03 (04 Feb 2024)-img/p2_chem_s1_q02_1.png}{68.5}{0}
\qfigdef{p2_chem_s1_q03_1}{Nurture SRG - Test 03 (04 Feb 2024)-img/p2_chem_s1_q03_1.png}{156.0}{0}
\qfigdef{p2_chem_s1_q04_1}{Nurture SRG - Test 03 (04 Feb 2024)-img/p2_chem_s1_q04_1.png}{159.4}{0}
\qfigdef{p2_chem_s1_q04_A}{Nurture SRG - Test 03 (04 Feb 2024)-img/p2_chem_s1_q04_A.png}{45.6}{0}
\qfigdef{p2_chem_s1_q04_B}{Nurture SRG - Test 03 (04 Feb 2024)-img/p2_chem_s1_q04_B.png}{48.9}{0}
\qfigdef{p2_chem_s1_q04_C}{Nurture SRG - Test 03 (04 Feb 2024)-img/p2_chem_s1_q04_C.png}{45.8}{0}
\qfigdef{p2_chem_s1_q04_D}{Nurture SRG - Test 03 (04 Feb 2024)-img/p2_chem_s1_q04_D.png}{51.9}{0}
\qfigdef{p2_chem_s2_q05_A}{Nurture SRG - Test 03 (04 Feb 2024)-img/p2_chem_s2_q05_A.png}{154.6}{0}
\qfigdef{p2_chem_s2_q05_B}{Nurture SRG - Test 03 (04 Feb 2024)-img/p2_chem_s2_q05_B.png}{206.3}{0}
\qfigdef{p2_chem_s2_q05_C}{Nurture SRG - Test 03 (04 Feb 2024)-img/p2_chem_s2_q05_C.png}{209.1}{0}
\qfigdef{p2_chem_s2_q05_D}{Nurture SRG - Test 03 (04 Feb 2024)-img/p2_chem_s2_q05_D.png}{208.9}{0}
\qfigdef{p2_chem_s2_q06_A}{Nurture SRG - Test 03 (04 Feb 2024)-img/p2_chem_s2_q06_A.png}{49.9}{0}
\qfigdef{p2_chem_s2_q06_B}{Nurture SRG - Test 03 (04 Feb 2024)-img/p2_chem_s2_q06_B.png}{26.2}{0}
\qfigdef{p2_chem_s2_q06_C}{Nurture SRG - Test 03 (04 Feb 2024)-img/p2_chem_s2_q06_C.png}{47.8}{0}
\qfigdef{p2_chem_s2_q06_D}{Nurture SRG - Test 03 (04 Feb 2024)-img/p2_chem_s2_q06_D.png}{51.1}{0}
\qfigdef{p2_chem_s2_q07_1}{Nurture SRG - Test 03 (04 Feb 2024)-img/p2_chem_s2_q07_1.png}{156.5}{0}
\qfigdef{p2_chem_s2_q07_A}{Nurture SRG - Test 03 (04 Feb 2024)-img/p2_chem_s2_q07_A.png}{81.1}{0}
\qfigdef{p2_chem_s2_q09_1}{Nurture SRG - Test 03 (04 Feb 2024)-img/p2_chem_s2_q09_1.png}{154.6}{0}
\qfigdef{p2_chem_s2_q10_A}{Nurture SRG - Test 03 (04 Feb 2024)-img/p2_chem_s2_q10_A.png}{196.2}{0}
\qfigdef{p2_chem_s2_q10_C}{Nurture SRG - Test 03 (04 Feb 2024)-img/p2_chem_s2_q10_C.png}{185.2}{0}
\qfigdef{p2_chem_s2_q11_1}{Nurture SRG - Test 03 (04 Feb 2024)-img/p2_chem_s2_q11_1.png}{186.4}{0}
\qfigdef{p2_chem_s2_q12_A}{Nurture SRG - Test 03 (04 Feb 2024)-img/p2_chem_s2_q12_A.png}{97.9}{0}
\qfigdef{p2_chem_s2_q12_B}{Nurture SRG - Test 03 (04 Feb 2024)-img/p2_chem_s2_q12_B.png}{107.5}{0}
\qfigdef{p2_chem_s2_q12_D}{Nurture SRG - Test 03 (04 Feb 2024)-img/p2_chem_s2_q12_D.png}{101.3}{0}
\qfigdef{p2_chem_s3_q02_1}{Nurture SRG - Test 03 (04 Feb 2024)-img/p2_chem_s3_q02_1.png}{283.8}{1}
\qfigdef{p2_chem_s3_q06_1}{Nurture SRG - Test 03 (04 Feb 2024)-img/p2_chem_s3_q06_1.png}{245.2}{0}
\qfigdef{p2_math_s1_q02_1}{Nurture SRG - Test 03 (04 Feb 2024)-img/p2_math_s1_q02_1.png}{73.2}{0}
\raggedcolumns
\makeatletter\@beginparpenalty=10000\makeatother
\newcolumntype{L}[1]{>{\raggedright\arraybackslash}p{#1}}
\begin{document}
\begin{center}{\Large\bfseries Nurture SRG --- Test 03 (04 Feb 2024)}\end{center}
\begin{multicols}{2}
\raggedcolumns
\papersec{1}
\partsec{1}{PHYSICS}
\sect{SECTION-I(i)}
\qn{1} Water in a clean aquarium forms a meniscus, as shown in the figure.\par
$T$ = surface tension of liquid\par
$h$ = difference of height between centre and the edge\par
$\rho$ = density of liquid\par
\qfig{p1_phy_s1_q01_1}
Walls of aquarium is large plane and contact angle is 0.
\begin{qoptions}
\opt{A}{$h\propto\dfrac T\rho$}
\opt{B}{$h\propto\left(\dfrac T\rho\right)^{1/2}$}
\opt{C}{$h\propto\left(\dfrac T\rho\right)^{2}$}
\opt{D}{$h\propto\left(\dfrac T\rho\right)^{3/2}$}
\end{qoptions}
\qn{2} A sealed cubical container with volume $V$ contains 20 particles each of mass $m$ but having various speeds. Two particles have speed $V$, three particles have speed $2V$, five particles have speed $3V$, four particles have speed $4V$, three particles have speed $5V$, two particles have speed $6V$ and one has speed $7V$. If the pressure exerted by these particles on the walls of container is $P$, then $PV$ is equal to
\begin{qoptions}
\opt{A}{$\dfrac{319}3mv^2$}
\opt{B}{$\dfrac83mv^2$}
\opt{C}{$\dfrac{51}3mv^2$}
\opt{D}{$\dfrac{20}3mv^2$}
\end{qoptions}
\qn{3} A ring of mass $m$ can freely slide on a smooth vertical rod. The ring symmetrically attached with two springs (as shown in the figure). Each of spring constant $k$. Each spring makes an angle $\theta$ with the horizontal. If the ring is slightly displaced vertically, its time period is:\par
\qfig{p1_phy_s1_q03_1}
\begin{qoptions}
\opt{A}{$2\pi\sqrt{\dfrac m{2k}}$}
\opt{B}{$2\pi\sqrt{\dfrac mk}$}
\opt{C}{$2\pi\sqrt{\dfrac{2m}{k\sin^2\theta}}$}
\opt{D}{$2\pi\sqrt{\dfrac m{2k\sin^2\theta}}$}
\end{qoptions}
\qn{4} How long will it take sound waves to travel to distance $\ell$ between the point A and B if the air temperature between them varies linearly from $T_1$ to $T_2$? The velocity sound propagation in air is equal to $v=\alpha\sqrt T$, where $\alpha$ is a constant.
\begin{qoptions}
\opt{A}{$t=\dfrac{\ell}{\alpha\left(\sqrt{T_1}+\sqrt{T_2}\right)}$}
\opt{B}{$t=\dfrac{\ell}{2\alpha\left(\sqrt{T_1}+\sqrt{T_2}\right)}$}
\opt{C}{$t=\dfrac{2\ell}{\alpha\left(\sqrt{T_1}+\sqrt{T_2}\right)}$}
\opt{D}{$t=\dfrac{\ell}{3\alpha\left(\sqrt{T_1}+\sqrt{T_2}\right)}$}
\end{qoptions}
\qn{5} Five sinusoidal waves have the same frequency 500 Hz but their amplitudes are in the ratio $2:1/2:1/2:1:1$ and their phase angles $0,\pi/6,\pi/3,\pi/2$ and $\pi$, respectively. The phase angle of resultant wave obtained by the superposition of these five waves is
\begin{qoptions}
\opt{A}{$30^\circ$}
\opt{B}{$45^\circ$}
\opt{C}{$60^\circ$}
\opt{D}{$90^\circ$}
\end{qoptions}
\qn{6} A stationary source is emitting sound at a fixed frequency $f_0$, which is reflected by two cars approaching the source (from same side). The difference between the frequencies of sound reflected from the cars is $1.2\%$ of $f_0$. What is the difference in the speed of the cars (in km per hour)? The cars are moving at constant speeds much smaller than the speed of sound which is $330\ \mathrm{ms^{-1}}$.
\begin{qoptions}
\opt{A}{6.128}
\opt{B}{7.128}
\opt{C}{8.128}
\opt{D}{9.128}
\end{qoptions}
\sect{SECTION-I(ii)}
\qn{7} Density of a liquid varies with depth as $\rho=\alpha h$. A small ball of density $\rho_0$ is released from the free surface of the liquid. Neglect viscous forces. Then
\begin{qoptions}
\opt{A}{The ball will execute SHM of amplitude $\dfrac{\rho_0}\alpha$}
\opt{B}{The mean position of the ball will be at a depth $\dfrac{\rho_0}\alpha$ from the free surface}
\opt{C}{The ball will sink to a maximum depth of $\dfrac{2\rho_0}\alpha$}
\opt{D}{The ball will sink to a maximum depth of $\dfrac{\rho_0}\alpha$}
\end{qoptions}
\qn{8} A source of sound and detector are moving as shown in figure at $t=0$. take velocity of sound wave to be $b$.\par
\qfig{p1_phy_s2_q08_1}
For the situation mark out the correct statement(s).
\begin{qoptions}
\opt{A}{The frequency received by the detector is always greater then $f_0$.}
\opt{B}{Initially, frequency received by the detector is greater than $f_0$, and then decreases with the time.}
\opt{C}{Frequency received by the detector is equal to $f_0$ at $t=d\cot\theta_0/(2v_0)$.}
\opt{D}{Frequency received by the detector can never be equal to $f_0$.}
\end{qoptions}
\qn{9} The disc of mass $m$ is connected to two springs $k_1$ and $k_2$ as shown, choose the correct option(s):\par
\qfig{p1_phy_s2_q09_1}
\begin{qoptions}
\opt{A}{$T=2\pi\sqrt{\dfrac{2m}{3\left(k_1+4k_2\right)}}$}
\opt{B}{$\omega=\sqrt{\dfrac{3\left(k_1+4k_2\right)}{2m}}$}
\opt{C}{$T=2\pi\sqrt{\dfrac{3m}{2\left(k_1+4k_2\right)}}$}
\opt{D}{$\omega=\sqrt{\dfrac{2\left(k_1+4k_2\right)}{3m}}$}
\end{qoptions}
\qn{10} Length of a string of density $\rho$ and Young's modulus $Y$ under tension is increased by $1/n$ times of its original length. If the velocity of transverse and longitudinal vibrations of the string is same, and the value of such velocity is $Y^a\rho^bn^c$, then:
\begin{qoptions}
\opt{A}{$a=-\dfrac12$}
\opt{B}{$b=-\dfrac12$}
\opt{C}{$c=-\dfrac13$}
\opt{D}{$c=-\dfrac14$}
\end{qoptions}
\qn{11} A metal wire length $L$, cross section area $A$, and Young's modulus $Y$ is stretched by a variable force $F$. $F$ is varying in such a ways that $F$ is always slightly greater than the elastic forces of resistance in the wire. When the elongation in the wire is $\ell$, up to this instant
\begin{qoptions}
\opt{A}{The work done by $F$ is $\dfrac{YA\ell^2}{2L}$}
\opt{B}{The work done by $F$ is $\dfrac{YA\ell^2}{L}$}
\opt{C}{The elastic potential energy stored in wire is $\dfrac{YA\ell^2}{2L}$}
\opt{D}{No energy is lost during elongation}
\end{qoptions}
\qn{12} Equal volumes of liquid are poured in the are poured in the three vessels A, B and C $(h_1<h_2<h_3)$. All the vessels have same base area. Select the correct alternatives. (Ignore atmospheric pressure)\par
\qfig{p1_phy_s2_q12_1}
\begin{qoptions}
\opt{A}{The force on the base will be maximum in vessel A}
\opt{B}{The force on the base will be maximum in vessel C}
\opt{C}{Net force exerted by the liquid on all the three vessels is equal}
\opt{D}{Net force exerted by the liquid in vessel A is maximum}
\end{qoptions}
\sect{SECTION-II}
\qn{1} A long metal rod of length $\ell$ and relative density $\sigma$ is held vertically with its lower end just touching the surface of water. The speed of the rod when it just sinks in water is given by (in m/s) (Take $\sigma=2$, $\ell=15$ cm)
\qn{2} A 1.2 m long metal wire is fixed securely at both ends to two solid supports so that the wire is initially horizontal. When a 29 g mass is attached from the midpoint of the wire, the midpoint is observed to move down by 20 mm. If the diameter of the wire is 0.1 mm, the Young's modulus for the wire material is $9.8\times10^x$. Find $x$.
\qn{3} If the square of angular frequency of small oscillation of block $m$ in the arrangement shown is equal to $\dfrac{3k_1k_3+\frac34k_1k_2+3k_2k_3}{\left(nk_3+\frac34k_2\right)m}$, find the value of $n$. Rod is massless. [Assume gravity to be absent]\par
\qfig{p1_phy_s3_q03_1}
\qn{4} A source is moving across a circle given by the equation $x^2+y^2=R^2$ with constant speed $v_s$ in clockwise sense. A detector is stationary at the point $(2R,0)$ w.r.t. the centre of the circle. The frequency emitted by the source is $f_s$. Find displacement between 2 positions of the source when the detector records the maximum and minimum frequencies. (given $R=2\sqrt3$m)
\qn{5} A stone is hung in air from a wire which is stretched over a sonometer. The bridges of the sonometer are 40 cm apart when the wire is in unison with a tuning fork of frequency 256 Hz. When the stone is completely immersed in water, the length between the bridges is 22 cm for re-establishing unison. The specific gravity of the material of the stone is (assume fundamental mode of vibration).
\qn{6} A light rod of length $L=2$ m is suspended horizontally from the ceiling by two wires A and B of equal lengths. The wire A is made of steel with the area of cross section $A_s=1\times10^{-5}\mathrm{m^2}$, while the wire B is made of brass of cross sectional area $A_b=2\times10^{-5}\mathrm{m^2}$. A weight $W$ is suspended at a distance $x$ from the wire A as shown in figure, Determine the value of $x$ in meters so that equal stresses are produced in each wire.\par
\qfig{p1_phy_s3_q06_1}
\partsec{2}{CHEMISTRY}
\sect{SECTION-I(i)}
\qn{1} Which of the following is INCORRECT statement ?
\begin{qoptions}
\opt{A}{\qfigi{p1_chem_s1_q01_A}}
\opt{B}{\qfigi{p1_chem_s1_q01_B}}
\opt{C}{\qfigi{p1_chem_s1_q01_C}}
\opt{D}{\qfigi{p1_chem_s1_q01_D}\par is more acidic than $\mathrm{CH_2{=}CH{-}CH_2{-}CH{=}CH_2}$}
\end{qoptions}
\qn{2} Match list I with list II and select the correct answer from the given below the lists :\par
{\footnotesize\setlength{\tabcolsep}{2pt}\begin{tabular}{|c|c|c|l|}
\hline
\multicolumn{2}{|c|}{List - I} & & List - II\\
\hline
P. & \qfigi{p1_chem_s1_q02_1} & i. & $(2R,3E)-3-$penten$-2-$ol\\
\hline
Q. & \qfigi{p1_chem_s1_q02_2} & ii. & $(2R,3Z)-3-$penten$-2-$ol\\
\hline
R. & \qfigi{p1_chem_s1_q02_3} & iii. & $(2S,3E)-3-$penten$-2-$ol\\
\hline
S. & \qfigi{p1_chem_s1_q02_4} & iv. & $(2S,3Z)-3-$penten$-2-$ol\\
\hline
\end{tabular}}
\begin{qoptions}
\opt{A}{P--i, Q--iii, R--iv, S--ii}
\opt{B}{P--i, Q--iii, R--ii, S--iv}
\opt{C}{P--iv, Q--iii, R--i, S--ii}
\opt{D}{P--iv, Q--iii, R--ii, S--i}
\end{qoptions}
\qn{3} During the following reaction, which intermediate is dominantly appeared ?\par
\qfig{p1_chem_s1_q03_1}
\begin{qoptions}
\opt{A}{\qfigi{p1_chem_s1_q03_A}}
\opt{B}{\qfigi{p1_chem_s1_q03_B}}
\opt{C}{\qfigi{p1_chem_s1_q03_C}}
\opt{D}{\qfigi{p1_chem_s1_q03_D}}
\end{qoptions}
\qn{4} The major product of the following reaction is\par
\qfig{p1_chem_s1_q04_1}
\begin{qoptions}
\opt{A}{\qfigi{p1_chem_s1_q04_A}}
\opt{B}{\qfigi{p1_chem_s1_q04_B}}
\opt{C}{\qfigi{p1_chem_s1_q04_C}}
\opt{D}{\qfigi{p1_chem_s1_q04_D}}
\end{qoptions}
\qn{5} 0.34 g of a hydrocarbon when heated with methyl magnesium bromide gives 112 mL of $\mathrm{CH_4}$ at NTP. Possible structure of the hydrocarbon is
\begin{qoptions}
\opt{A}{$\mathrm{CH_3CH_2CH_2C{\equiv}CH}$}
\opt{B}{$\mathrm{CH_3CH_2C{\equiv}C{-}CH_3}$}
\opt{C}{$\mathrm{CH_3CH_2{-}CH(CH_3){-}C{\equiv}CH}$}
\opt{D}{$\mathrm{CH_3CH_2CH_2CH_2C{\equiv}CH}$}
\end{qoptions}
\qn{6} \qfig{p1_chem_s1_q06_1}
\begin{qoptions}
\opt{A}{\qfigi{p1_chem_s1_q06_A}}
\opt{B}{\qfigi{p1_chem_s1_q06_B}}
\opt{C}{\qfigi{p1_chem_s1_q06_C}}
\opt{D}{\qfigi{p1_chem_s1_q06_D}}
\end{qoptions}
\sect{SECTION-I(ii)}
\qn{7} Which of the following compound will not give Friedel-Crafts reaction ?
\begin{qoptions}
\opt{A}{\qfigi{p1_chem_s2_q07_A}}
\opt{B}{\qfigi{p1_chem_s2_q07_B}}
\opt{C}{\qfigi{p1_chem_s2_q07_C}}
\opt{D}{\qfigi{p1_chem_s2_q07_D}}
\end{qoptions}
\qn{8} $\mathrm{H{-}C{\equiv}CLi}$ + $\mathrm{Br{-}(CH_2)_8{-}Br}\to A$ $\xrightarrow{\mathrm{CH_3Li}}[B]\xrightarrow{A}\underset{\text{(possible)}}{C}$\par
In this reaction, sequences
\begin{qoptions}
\opt{A}{$A$ is $\mathrm{H{-}C{\equiv}C{-}(CH_2)_8{-}Br}$}
\opt{B}{$B$ is $\mathrm{H{-}C{\equiv}C{-}H}$ and $C$ is $\mathrm{Li\,C{\equiv}C{-}(CH_2)_8{-}Br}$}
\opt{C}{$B$ is $\mathrm{Li\,C{\equiv}C{-}(CH_2)_8{-}Br}$ and $C$ is\par \qfigi{p1_chem_s2_q08_1}}
\opt{D}{$B$ is $\mathrm{Li\,C{\equiv}C{-}(CH)_8{-}Br}$ and $C$ is $\mathrm{Br(CH_2)_8}$ $\mathrm{C{\equiv}C(CH_2)_8Br}$}
\end{qoptions}
\qn{9} Which one of the following is/are \textbf{INCORRECT} ?
\begin{qoptions}
\opt{A}{Fluorapatite is much harder than hydroxyapatite}
\opt{B}{Max. allowed limit of Pb in water is 50 ppm}
\opt{C}{Excess of sulphate in water causes blue baby syndrome}
\opt{D}{Excess of nitrate in water causes constipation}
\end{qoptions}
\qn{10} Which of the following reactions involve free radical as intermediate ?
\begin{qoptions}
\opt{A}{\qfigi{p1_chem_s2_q10_A}}
\opt{B}{\qfigi{p1_chem_s2_q10_B}}
\opt{C}{\qfigi{p1_chem_s2_q10_C}}
\opt{D}{\qfigi{p1_chem_s2_q10_D}}
\end{qoptions}
\qn{11} Which of the following is/are correct statements ?\par
\qfig{p1_chem_s2_q11_1}
\begin{qoptions}
\opt{A}{$Q$ is \qfigi{p1_chem_s2_q11_2}}
\opt{B}{Stereomers of product $R$ involves internal as well as external compensation of optical rotation}
\opt{C}{Single enantiomerically pure form of $Q$ is formed}
\opt{D}{Enantiomers as well as diastereomers are formed as $R$}
\end{qoptions}
\qn{12} \qfig{p1_chem_s2_q12_1}
\begin{qoptions}
\opt{A}{$A$ is \qfigi{p1_chem_s2_q12_2}}
\opt{B}{$B$ is \qfigi{p1_chem_s2_q12_3}}
\opt{C}{$C$ has 3 stereogenic centres}
\opt{D}{$B$ has the same number of stereogenic centres as $C$}
\end{qoptions}
\sect{SECTION-II}
\qn{1} Find the total number of possible alkyl chloride (R-Cl) in the following reaction.\par
\qfig{p1_chem_s3_q01_1}
\qn{2} What volume (lit) of ethane is formed from 38 g of sodium propionate by fusion with sode lime ? (At NTP, 1mole of ideal gas occupies 22.4 lit volume)
\qn{3} Consider the following reactions\par
\qfigi{p1_chem_s3_q03_1} Total number of monochlorinated product = X (Excluding stereoisomers)\par
\qfigi{p1_chem_s3_q03_2} Total number of monochlorinated product = Y (Excluding stereoisomers)\par
Then, the value of $(X+Y)$ is
\qn{4} Identify number of compounds which react faster than benzene in ArSE reaction ?\par
\qfig{p1_chem_s3_q04_1}
\qn{5} How many compounds are soluble in aqueous NaOH ?\par
\qfig{p1_chem_s3_q05_1}
\qn{6} Among the following, the number of aromatic compound(s) is(are)\par
\qfig{p1_chem_s3_q06_1}
\partsec{3}{MATHEMATICS}
\sect{SECTION-I(i)}
\qn{1} If $Z_r$; $r=1,2,3,\dots,50$ are the roots of the equation $\displaystyle\sum_{r=0}^{50}(Z)^r=0$, then the value of $\displaystyle\sum_{r=1}^{50}\dfrac1{Z_r-1}$ is
\begin{qoptions}
\opt{A}{$-85$}
\opt{B}{$-25$}
\opt{C}{$25$}
\opt{D}{$75$}
\end{qoptions}
\qn{2} If $z$ \& $w$ are two complex numbers simultaneously satisfying the equations, $z^3+w^5=0$ and $z^2\overline w^4=1$, then:
\begin{qoptions}
\opt{A}{$z$ and $w$ both are purely real}
\opt{B}{$z$ is purely real and $w$ is purely imaginary}
\opt{C}{$w$ is purely real and $z$ is purely imaginary}
\opt{D}{$z$ and $w$ both are imaginary}
\end{qoptions}
\qn{3} Any ordinate $MP$ of an ellipse $\dfrac{x^2}{25}+\dfrac{y^2}9=1$ meets the auxiliary circle in $Q$, Then locus of point of intersection of normals at $P$ and $Q$ to the respective curves is
\begin{qoptions}
\opt{A}{$x^2+y^2=8$}
\opt{B}{$x^2+y^2=34$}
\opt{C}{$x^2+y^2=64$}
\opt{D}{$x^2+y^2=15$}
\end{qoptions}
\qn{4} Find the range of parameter `$a$' for which a unique circle will pass through the point of intersection of the rectangular hyperbola $x^2-y^2=a^2$ and parabola $y=x^2$
\begin{qoptions}
\opt{A}{$a\in(-1,1)$}
\opt{B}{$a\in(0,1)$}
\opt{C}{$a\in\left(\dfrac{-1}2,\dfrac12\right)$}
\opt{D}{$a\in\left(\dfrac{-1}2,\dfrac14\right)$}
\end{qoptions}
\qn{5} If $z_1,z_2$ and $z_3$ are the vertices of $\Delta ABC$, which is not right angled triangle taken in anti-clock wise direction and $z_0$ is the circumcentre, then $\left(\dfrac{z_0-z_1}{z_0-z_2}\right)\dfrac{\sin2A}{\sin2B}+\left(\dfrac{z_0-z_3}{z_0-z_2}\right)\dfrac{\sin2C}{\sin2B}$ is equal to
\begin{qoptions}
\opt{A}{0}
\opt{B}{1}
\opt{C}{$-1$}
\opt{D}{2}
\end{qoptions}
\qn{6} Let $P$ denotes a complex number $z=r(\cos\theta+i\sin\theta)$ on the Argand's plane, and $Q$ denotes a complex number $\sqrt{2|z|^2}\left(\cos\left(\theta+\dfrac\pi4\right)+i\sin\left(\theta+\dfrac\pi4\right)\right)$. If `O' is the origin, then $\Delta OPQ$ is
\begin{qoptions}
\opt{A}{Isosceles but not right angled}
\opt{B}{Right angled but not isosceles}
\opt{C}{Right isosceles}
\opt{D}{Equilateral}
\end{qoptions}
\sect{SECTION-I(ii)}
\qn{7} If $z_1=5+12i$ and $|z_2|=4$ then
\begin{qoptions}
\opt{A}{Maximum $\left(|z_1+iz_2|\right)=17$}
\opt{B}{Minimum $\left(|z_1+(1+i|z_2|\,|\right)=13-9\sqrt2$}
\opt{C}{Minimum $\left|\dfrac{z_1}{z_2+\frac4{z_2}}\right|=\dfrac{13}4$}
\opt{D}{Maximum $\left|\dfrac{z_1}{z_2+\frac4{z_2}}\right|=\dfrac{13}4$}
\end{qoptions}
\qn{8} Line $OQ$ is angle bisector of angle $O$ of right angle triangle $OPR$, right angle at $P$. Point $Q$ is such that $ORQP$ is concyclic. If point $O$ is origin and points $P,Q,R$ are represented by the complex numbers $Z_3$, $Z_2$, $Z_1$ respectively. If $\dfrac{Z_2^2}{Z_3Z_2}=\dfrac32$ then ($R$ is circum radius of $\Delta OPR$)
\begin{qoptions}
\opt{A}{Angles of $\Delta OPR$ are $\dfrac\pi6,\dfrac\pi3,\dfrac\pi2$}
\opt{B}{Angles of $\Delta OPR$ are $\dfrac\pi4,\dfrac\pi4,\dfrac\pi2$}
\opt{C}{Area of $\Delta OPR$ is $2\sqrt2R^2$}
\opt{D}{Area of $\Delta OPR$ is $\dfrac{\sqrt3}2R^2$}
\end{qoptions}
\qn{9} Let $z_1,z_2,z_3$ be the coordinate of verticles A, B, C of the triangle if $W_1=z_1-z_2$, $W_2=z_3-z_2$ and $\angle A=90^\circ$
\begin{qoptions}
\opt{A}{$\mathrm{Re}\left(\dfrac{W_2}{W_1}\right)=1$}
\opt{B}{$\mathrm{Re}\left(W_1\overline{W_2}\right)=|W_1|^2$}
\opt{C}{$\mathrm{Re}\left(W_1\overline{W_2}\right)=|W_2|^2$}
\opt{D}{$\mathrm{Re}\left(\dfrac{W_1}{W_2}\right)=1$}
\end{qoptions}
\qn{10} Let $S=\{1,2,3,\dots,20\}$ be the set of all positive integers from 1 to 20. Suppose that $N$ is the smallest positive integer such that exactly eighteen numbers from $S$ are factors of $N$, and the only two numbers from $S$ that are not factors of $N$ are consecutive integers. If $K$ is the total number of positive integral divisors of $N$, then $K$ is greater than
\begin{qoptions}
\opt{A}{200}
\opt{B}{300}
\opt{C}{400}
\opt{D}{500}
\end{qoptions}
\qn{11} Let a variable circle touches a line, which is tangent to the parabola $y^2=4x$ at $(4,4)$ and variable circle passes through focus of the parabola $y^2=4x$. If locus of its centre is conic `C' then
\begin{qoptions}
\opt{A}{Directrix of the conic `C' is $x-2y+4=0$}
\opt{B}{Focus of the conic `C' is $(4,0)$}
\opt{C}{Latus rectum of conic `C' is $2\sqrt5$}
\opt{D}{Axis of the conic `C' is $x+2y-1=0$}
\end{qoptions}
\qn{12} Let $P(t^2,2t)$ be the point on the parabola $y^2=4x$ other than its vertex $A$ whose abscissa and ordinate are same. A circle $C$ is drawn with focal radius of the point $P$ as diameter, then which of the following hold(s) good ?
\begin{qoptions}
\opt{A}{Circle $C$ touches the directrix of the parabola.}
\opt{B}{Circle $C$ touches the tangent at $A$ of the parabola.}
\opt{C}{Length of the intercept made by the circle $C$ on the normal drawn at $P$ of the parabola is $\sqrt5$}
\opt{D}{Radius of the circle $C$ is 5.}
\end{qoptions}
\sect{SECTION-II}
\qn{1} Suppose we are constructing 12 lettered word with four X's four Y's and four Z's. Let $p$ denotes the number of 12 lettered words which have no X's in the first four letters, and $q$ denotes the number of 12 lettered words which have no X's in the first four letters, no Y's in the middle four letters and no Z's in the last four letters. Find $(p+q)$.
\qn{2} $PQ$ is a variable chord of the parabola $y^2=4x$ such that its projection with a line with gradient unity has a magnitude equal to 10. If the locus of the middle point of $PQ$ has the equation $\left(y^2-4x\right)(y+2)^2+\lambda=0$ where $\lambda\in N$, then find $\lambda$.
\qn{3} Find the minimum value of the expression $E=|z|^2+|z-3|^2+|z-6i|^2$ (where $z=x+iy$, $x,y\in R$)
\qn{4} In a paper of mathematics there are 5 questions, such that the sum of the marks is 30 and the marks for any question is not less than 2 and not more than 8 (only integer marks are awarded). If the number of ways in which the marks can be awarded is $ABC$, where $A,B,C\in N$ and $A,B,C\le9$. Then value of $A+B-C$ is \underline{\hspace{1cm}}
\qn{5} Let the chord $x\cos\alpha+y\sin\alpha=p$ of the hyperbola $\dfrac{x^2}{16}-\dfrac{y^2}{18}=1$ subtends a right angle at the centre. Let diameter of the circle, concentric with the hyperbola, to which the given chord is a tangent is $d$. Then $\dfrac d4=$?
\qn{6} There are 20 observations $x_1,x_2,\dots,x_{20}$. If mean of $(x_i+i)\ \forall\ i=1,2,3,\dots20$ is 15.5 and $\displaystyle\sum_{i=1}^{20}(x_i-5)^2=100$. If variance of $(x_i-2)\forall\ i=1,2,3\dots,20$ is $\lambda$, then $\dfrac\lambda4$ is equal to
\papersec{2}
\partsec{1}{PHYSICS}
\sect{SECTION-I(i)}
\qn{1} A thermally insulated vessel contains some water at $0^\circ$C. The vessel is connected to a vacuum pump to pump out water vapour. This results in some water getting frozen. It is given Latent heat of vaporization of water at $0^\circ$C $=21\times10^5$ J/kg and latent heat of freezing of water $=3.36\times10^5$ J/kg. The maximum percentage amount of water that will be solidified in this manner will be :
\begin{qoptions}
\opt{A}{$86.2\%$}
\opt{B}{$33.6\%$}
\opt{C}{$21\%$}
\opt{D}{$24.36\%$}
\end{qoptions}
\qn{2} A cylindrical tube of cross-sectional area $A$ has two air tight frictionless pistons at its two ends. The pistons are tied with a straight piece of metallic wire. The tube contains a gas at atmospheric pressure $P_0$ and temperature $T_0$, If temperature of the gas is doubled then the tension in the wire is\par
\qfig{p2_phy_s1_q02_1}
\begin{qoptions}
\opt{A}{$4P_0A$}
\opt{B}{$P_0A/2$}
\opt{C}{$P_0A$}
\opt{D}{$2P_0A$}
\end{qoptions}
\qn{3} A bar in figure below is made from a single piece of material having Young's modulus $Y$. It consists of two segments of equal length $\dfrac{L_0}2$ and cross-sectional area $A$ and $2A$. The elastic potential energy stored under the action of an axial force $F$ is\par
\qfig{p2_phy_s1_q03_1}
\begin{qoptions}
\opt{A}{$\dfrac{3}{4AY}F^2L_0$}
\opt{B}{$\dfrac{3}{AY}F^2L_0$}
\opt{C}{$\dfrac{12}{AY}F^2L_0$}
\opt{D}{$\dfrac{3}{8AY}F^2L_0$}
\end{qoptions}
\qn{4} A near surface of earth satellite is in the shape of a sphere of radius $r$. It encounters cosmic dust in its path. The viscous force experienced by the satellite follows Stoke's law. The coefficient of viscosity is $\eta$. Mass and radius of the earth are $M$ and $R$ respectively. The power of the rocket engine that must be put on to keep the satellite moving as usual (if $\eta GM=5\times10^{10}$ in MKS and $\dfrac rR=\dfrac{10^{-5}}\pi$) is (orbital speed of satellite is $\sqrt{\dfrac{GM}R}$)
\begin{qoptions}
\opt{A}{$3\times10^6$ Watt}
\opt{B}{$4\times10^6$ Watt}
\opt{C}{$5\times10^6$ Watt}
\opt{D}{$6\times10^6$ Watt}
\end{qoptions}
\sect{SECTION-I(ii)}
\qn{5} A paraboloid shape container is placed at origin as shown in the figure, filled with liquid up to a height 2m. Equation of parabola is given as $x^2=2y$. A small opening is made at origin of area $10\ \mathrm{cm^2}$, then $(\pi=\sqrt{10})$ $(g=10\ \mathrm{m/s^2})$\par
\qfig{p2_phy_s2_q05_1}
\begin{qoptions}
\opt{A}{Speed of efflux when height of liquid is 0.8 m is 2 m/s}
\opt{B}{Speed of efflux when height of liquid is 0.2 m is 2 m/s}
\opt{C}{Time to empty the container is $\dfrac{20}{27}$hr}
\opt{D}{Time to empty the container is $\dfrac{27}{20}$hr}
\end{qoptions}
\qn{6} Let $v_{av}$, $v_{rms}$ and $v_p$ respectively denote mean speed, root mean square speed and the most probable speed of the molecule in an ideal monoatomic gas at absolute temperature $T$. The mass of a molecule is $m$ then :
\begin{qoptions}
\opt{A}{No molecule can have speed greater than $\sqrt2\,v_{rms}$}
\opt{B}{No molecule can have speed less than $\dfrac{v_p}{\sqrt2}$}
\opt{C}{$v_p<v_{av}<v_{rms}$}
\opt{D}{The average kinetic energy of a molecule is $\dfrac34mv_p^2$}
\end{qoptions}
\qn{7} The bottom of the cylinder contains 2 kg of water at $100^\circ$C. The cylinder has a radius $\left(\dfrac{10}{\sqrt\pi}\right)$cm. The massless piston sits on the surface of water. An electric heater at the base of cylinder transfers energy into the water at the rate of 100 W. Assume the cylinder is of sufficient length so that piston never reaches the top of the cylinder. (Take $R=\dfrac{25}3$, $L_v=2.25\times10^6$J/kg and $S_{steam}=2000$ J/kg-$^\circ$C, $P_0=10^5\ \mathrm{N/m^2}$) Assume steam as an ideal gas.\par
\qfig{p2_phy_s2_q07_1}
\begin{qoptions}
\opt{A}{The rate of rise of piston when water starts boiling into steam is 7.67 mm/sec.}
\opt{B}{The rate of rise of piston when water starts boiling into steam is 1.2 mm/sec.}
\opt{C}{The rate of rise of piston when water has completely converted into steam is 12.6 mm/sec.}
\opt{D}{The rate of rise of piston when water has completely converted into steam is 23.1 mm/sec.}
\end{qoptions}
\qn{8} A travelling wave pulse on a string of linear mass density $\dfrac{3\mathrm g}{\mathrm{cm}}$ is given by $y=2.00e^{-\left(4x^2+324t^2+72xt+3\right)}$ where $x$ and $y$ are in meter and $t$ in second then choose the correct options.
\begin{qoptions}
\opt{A}{Wave pulse is travelling along positive x-axis with speed 9 m/s.}
\opt{B}{Wave pulse in travelling along negative x-axis with speed 9 m/s.}
\opt{C}{Tension in the string is 24.3 N}
\opt{D}{Amplitude of wave pulse is 2.00 m}
\end{qoptions}
\qn{9} A particle performing SHM is found at its equilibrium position at $t=1$ sec and it's velocity (in positive direction) is 0.25 m/s at $t=2$ sec. If the time period of oscillation is 6 sec. Then -
\begin{qoptions}
\opt{A}{Amplitude of oscillation is $\dfrac3{2\pi}$m}
\opt{B}{Speed of particle at 6 sec is 0.5 m/s}
\opt{C}{Speed of particle at 6 sec is 0.25 m/s}
\opt{D}{Acceleration of particle at 6 sec is zero.}
\end{qoptions}
\qn{10} For a body executing SHM with amplitude $A$, time period $T$, maximum velocity $V_{max}$ and phase constant zero, (At $t=0$ body at mean position). Which of the following statements are correct for $0\le t\le\dfrac T4$ ($y$ is displacement from mean position) ?
\begin{qoptions}
\opt{A}{At $y=(A/2)$, $v>(v_{max}/2)$}
\opt{B}{For $v=(v_{max}/2)$, $y>(A/2)$}
\opt{C}{For $t=(T/8)$, $y>(A/2)$}
\opt{D}{For $y=(A/2)$, $t<(T/8)$}
\end{qoptions}
\qn{11} A cubical container of side length `L' is half filled with liquid of density $\rho$ and placed on the floor of an elevator accelerating downward with acceleration `g/2'. Now the container is given a horizontal acceleration parallel to one of the horizontal edge then ($P_0$ = atmospheric pressure)
\begin{qoptions}
\opt{A}{Pressure at centre of container will remain unchanged for any value of horizontal acceleration}
\opt{B}{The horizontal acceleration for which liquid over flow takes place is $\dfrac g2$}
\opt{C}{The maximum gauge pressure at the bottom surface when liquid is about to over flow is $P_0+\dfrac{\rho gL}2$}
\opt{D}{The minimum pressure at the bottom surface when liquid is about to over flow is $P_0$}
\end{qoptions}
\qn{12} Consider a situation in which there is a Newtonian fluid between two long coaxial uniformly moving cylinders of radius $R$ and $2R$.\par
\qfig{p2_phy_s2_q12_1}
\begin{qoptions}
\opt{A}{Force per unit length required to keep the inner cylinder moving is $4\pi\eta v_0$}
\opt{B}{Force per unit length required to keep the inner cylinder moving is $\dfrac{6\pi\eta v_0}{\ell n2}$}
\opt{C}{Distance of the point from the axis at which speed of fluid is zero is $2^{1/3}R$}
\opt{D}{Distance of the point from the axis at which speed of fluid is zero is $\sqrt2R$}
\end{qoptions}
\sect{SECTION-II}
\qn{1} One cylindrical vessel of radius $R_1$ is fixed inside another cylindrical vessel of radius $R_2$, as shown in the figure. In the bottom of the small vessel, there is a small hole and a wooden cylinder of radius $r$ and height $h=21$ cm is inserted. The wooden cylinder can only move vertically without friction. Water is poured into the small vessel to a height of $a=30$ cm, and oil is poured into the large vessel to the same level and the wooden cylinder is in equilibrium. Given that the water density is $\rho_w=1000\mathrm{kg/m^3}$, the oil density is $\rho_0=790\mathrm{kg/m^3}$ and the wooden cylinder density is $\rho=600\mathrm{kg/m^3}$.\par
\qfig{p2_phy_s3_q01_1}
Find the length of wooden cylinder inside water at equilibrium in cm. Water remains in smaller cylinder.
\qn{2} A string fixed at both ends with tension 100N is vibrating in lowest mode of vibration for which a point at quarter of its length from one end is a point of maximum displacement. The frequency of vibration in this mode is 30 Hz. The frequency (in Hz) emitted when it vibrates in the next mode with new tension 25N such that this point is again a point of maximum displacement is
\qn{3} The position of a particle of mass 1 kg with respect to origin varies according to relation $x=3\sin(100t)+8\cos^250t$, where $x$ is in meter and $t$ in sec. Find out total energy of particle in (kJ) \{Assume PE at mean position is zero\}
\qn{4} Calculate the sound intensity level in decibel for a sound wave travelling in air at $0^\circ$C and having a pressure amplitude of 0.5 Pa. Assume density of air = 1.25 kg/m$^3$, velocity of sound = 330 m/s, $\log_{10}(3.3)=0.52$
\qn{5} A 5.0 long vertical pipe is filled completely with a liquid. A small hole is drilled at the base of pipe through which liquid starts leaking out. This pipe resonates with a tuning fork. The first time resonance occur when height of water column is 4.4 m and second time when height is 2.8 m respectively. If radius of pipe is found to be $\left(\dfrac1x\right)$m, find $x$
\qn{6} A particle of mass 1.5 kg moves along the x-axis in force field, Its potential energy is given by $V(x)=2x^3-9x^2+12x$, where all quantities are written in S.I units. There is an equilibrium point on the x-axis at which if displaced slightly it performs SHM. If the angular frequency of oscillation is $x$ rad/s. Find $x$.
\partsec{2}{CHEMISTRY}
\sect{SECTION-I(i)}
\qn{1} Which of the following compound will lose optical activity after the reaction?
\begin{qoptions}
\opt{A}{\qfigi{p2_chem_s1_q01_A}}
\opt{B}{\qfigi{p2_chem_s1_q01_B}}
\opt{C}{\qfigi{p2_chem_s1_q01_C}}
\opt{D}{\qfigi{p2_chem_s1_q01_D}}
\end{qoptions}
\qn{2} The INCORRECT statement Regarding Rearrangement in the given carbocation?\par
\qfig{p2_chem_s1_q02_1}
\begin{qoptions}
\opt{A}{It rearranges to most stable carbocation by ring Expansion}
\opt{B}{Its most stable form obtained after rearrangement has a resonance structure with completed valence shells of all the atoms.}
\opt{C}{To its most stable form obtained after rearrangement, oxygen decreases stability by inductive effect but increases stability by resonance.}
\opt{D}{$\mathrm{CH_3}$-shift is the major mode of Rearrangement.}
\end{qoptions}
\qn{3} Among the following compounds;\par
\qfig{p2_chem_s1_q03_1}
The correct decreasing order of Enol-content is;
\begin{qoptions}
\opt{A}{I, II, III}
\opt{B}{III, I, II}
\opt{C}{II, I, III}
\opt{D}{II, III, I}
\end{qoptions}
\qn{4} \qfig{p2_chem_s1_q04_1}
\begin{qoptions}
\opt{A}{\qfigi{p2_chem_s1_q04_A}}
\opt{B}{\qfigi{p2_chem_s1_q04_B}}
\opt{C}{\qfigi{p2_chem_s1_q04_C}}
\opt{D}{\qfigi{p2_chem_s1_q04_D}}
\end{qoptions}
\sect{SECTION-I(ii)}
\qn{5} Select the option which contains the CORRECT order against the parameter mentioned in the parenthesis
\begin{qoptions}
\opt{A}{\qfigi{p2_chem_s2_q05_A}}
\opt{B}{\qfigi{p2_chem_s2_q05_B}}
\opt{C}{\qfigi{p2_chem_s2_q05_C}}
\opt{D}{\qfigi{p2_chem_s2_q05_D}}
\end{qoptions}
\qn{6} Which of the following give only meso compound on catalytic reduction with Raney $\mathrm{Ni/H_2}$?
\begin{qoptions}
\opt{A}{\qfigi{p2_chem_s2_q06_A}}
\opt{B}{\qfigi{p2_chem_s2_q06_B}}
\opt{C}{\qfigi{p2_chem_s2_q06_C}}
\opt{D}{\qfigi{p2_chem_s2_q06_D}}
\end{qoptions}
\qn{7} Consider the following energy profile diagram for the given elimination reaction ($E_1$) of a $3^\circ$ alkyl halide $R_1$, $R_2$ and $R_3$ containing different halogen atom with same alkyl group.\par
\qfig{p2_chem_s2_q07_1}
Two possible products are obtained as $P_1$ and $P_2$. Identify the \textbf{CORRECT} statement(s)
\begin{qoptions}
\opt{A}{\qfigi{p2_chem_s2_q07_A}}
\opt{B}{$P_1$ is 2-methylbut-1-ene}
\opt{C}{$P_2$ has 9 hyperconjugable hydrogens}
\opt{D}{$P_1$ on reaction with $\mathrm{Ni/D_2}$ will give a racemic mixture}
\end{qoptions}
\qn{8} Select the correct options given below;
\begin{qoptions}
\opt{A}{NO is formed from $\mathrm{N_2}$ and $\mathrm{O_2}$, when lightening occurs at high altitude.}
\opt{B}{NO is formed when $\mathrm{N_2}$ is oxidised by $\mathrm{O_3}$}
\opt{C}{NO is oxidised to $\mathrm{NO_2}$ by $\mathrm{O_2}$ or by $\mathrm{O_3}$ in the stratosphere.}
\opt{D}{$\mathrm{N_2O}$ is formed by the reaction $\mathrm{N_2}$ with $\mathrm{O_3}$ or with atomic oxygen in the upper atmosphere and also by microbiological process.}
\end{qoptions}
\qn{9} \qfig{p2_chem_s2_q09_1}
Among the following which Statements are \textbf{CORRECT}?
\begin{qoptions}
\opt{A}{Both N of pyrimidine are of same basic strength}
\opt{B}{In imidazole, protonation takes place on N---3}
\opt{C}{Purine has 3 basic Nitrogen sites}
\opt{D}{Purine, Imidazole and pyrimidine all are aromatic}
\end{qoptions}
\qn{10} Which of these statement are correct?
\begin{qoptions}
\opt{A}{\qfigi{p2_chem_s2_q10_A}}
\opt{B}{Acetone behaves as electrophile as well as nucleophile}
\opt{C}{\qfigi{p2_chem_s2_q10_C}}
\opt{D}{$\mathrm{RO^{\ominus}}$ is less basic than $\mathrm{RS^{\ominus}}$}
\end{qoptions}
\qn{11} \qfig{p2_chem_s2_q11_1}
Relation between (B) and (C) is:
\begin{qoptions}
\opt{A}{Enantiomers}
\opt{B}{Diastereomers}
\opt{C}{Geometrical isomers}
\opt{D}{Homomers}
\end{qoptions}
\qn{12} One mole of an unknown organic compound (A) when treated with excess of $\mathrm{CH_3MgBr}$ liberates three moles of $\mathrm{CH_4}$ gas. When reduced with HI and red phosphorus, compound (A) gives n- butane. The possible structure(s) of A is(are):
\begin{qoptions}
\opt{A}{\qfigi{p2_chem_s2_q12_A}}
\opt{B}{\qfigi{p2_chem_s2_q12_B}}
\opt{C}{$\mathrm{HOCH_2{-}C{\equiv}C{-}CH_2OH}$}
\opt{D}{\qfigi{p2_chem_s2_q12_D}}
\end{qoptions}
\sect{SECTION-II}
\qn{1} How many monocarboxylic acids (including stereoisomers) would give Methylcyclopropane on sodalime decarboxylation?
\qfullbegin
\qn{2} How many compounds from following exhibit d-orbital resonance.\par
\qfig{p2_chem_s3_q02_1}
\qfullend
\qn{3} Number of diastereomers of $\mathrm{C_2(F)(C\ell)(Br)(I)}$ is $\mathbf x$,\par
Number of isomeric alcohols of $\mathrm{C_4H_{10}O}$ is $\mathbf y$ and\par
Number of isomeric esters of $\mathrm{C_4H_8O_2}$ is $\mathbf z$.\par
Calculate the value of $(\mathbf x+\mathbf y+\mathbf z)$.
\qn{4} Calculate total number of alkene products (including stereoisomers) formed when 3-chloro-3methyl heptane react with alcoholic KOH and heat.
\qn{5} Dichlorination of isopentane (excluding stereoisomers) gives $x$ structural isomers. The value of $x$ is
\qn{6} How many compounds are more reactive than benzene towards Nitration?\par
\qfig{p2_chem_s3_q06_1}
\partsec{3}{MATHEMATICS}
\sect{SECTION-I(i)}
\qn{1} If $\dfrac1{a+\omega}+\dfrac1{b+\omega}+\dfrac1{c+\omega}+\dfrac1{d+\omega}=\dfrac1\omega$, where $a,b,c,d\in R$ and $\omega$ is a cube root of unity then $\sum\dfrac1{a^2-a+1}$ is equal to
\begin{qoptions}
\opt{A}{1}
\opt{B}{2}
\opt{C}{3}
\opt{D}{$\dfrac13$}
\end{qoptions}
\qn{2} The number of ways in which all the letters of the word `BIGBOSS' can be placed in the squares of the figure shown, so that no row remains empty is equal to\par
\qfig{p2_math_s1_q02_1}
\begin{qoptions}
\opt{A}{95760}
\opt{B}{94640}
\opt{C}{97840}
\opt{D}{82120}
\end{qoptions}
\qn{3} The locus of the middle points of chords of the parabola $y^2=4x$, which are of constant length `$2\ell$' is :
\begin{qoptions}
\opt{A}{$\left(4x+y^2\right)\left(y^2-4\right)=4\ell^2$}
\opt{B}{$\left(4y+x^2\right)\left(x^2-4\right)=4\ell^2$}
\opt{C}{$\left(4y-x^2\right)\left(x^2+4\right)=4\ell^2$}
\opt{D}{$\left(4x-y^2\right)\left(y^2+4\right)=4\ell^2$}
\end{qoptions}
\qn{4} Point `O' is the centre of the ellipse with major axis AB and minor axis CD. Point F is one focus of the ellipse. If $OF=6$ and the diameter of the inscribed circle of triangle OCF is 2, then the value of $\sqrt{(AB)^2-(CD)^2}$
\begin{qoptions}
\opt{A}{36}
\opt{B}{6}
\opt{C}{44}
\opt{D}{12}
\end{qoptions}
\sect{SECTION-I(ii)}
\qn{5} The locus of incentre of $\Delta PF_1F_2$ where $P$ is a variable point lying on ellipse $\dfrac{x^2}{a^2}+\dfrac{y^2}{b^2}=1$ with $F_1,F_2$ are foci and eccentricity $e$ is a conic whose :
\begin{qoptions}
\opt{A}{Equation is $\dfrac{x^2}{a^2e^2}+\dfrac{y^2(1+e)^2}{b^2e^2}=1$}
\opt{B}{Equation is $\dfrac{x^2}{a^2e^2}+\dfrac{y^2(1-e)^2}{b^2e^2}=1$}
\opt{C}{Eccentricity is $\sqrt{\dfrac{2e}{1+e}}$}
\opt{D}{Eccentricity is $\sqrt{\dfrac{2e}{1-e}}$}
\end{qoptions}
\qn{6} Given ellipse $\dfrac{x^2}{16}+\dfrac{y^2}7=1$ and the hyperbola $\dfrac{x^2}{144}-\dfrac{y^2}{81}=\dfrac1{25}$, if the ordinate of one of the points of intersection is produced to cut asymptote at $P$, then which of the following is/are true ?
\begin{qoptions}
\opt{A}{They have the same foci}
\opt{B}{Square of the ordinate of point of intersection is $\dfrac{63}{25}$.}
\opt{C}{Sum of the square of coordinates of $P$ is 16}
\opt{D}{$P$ lies on the auxiliary circle formed by ellipse.}
\end{qoptions}
\qn{7} Let $z_1,z_2$ and $z_3$ be three distinct complex numbers, satisfying $|z_1|=|z_2|=|z_3|=1$. Which of the following is/are true :
\begin{qoptions}
\opt{A}{If $\arg\left(\dfrac{z_1}{z_2}\right)=\dfrac\pi2$ then $\arg\left(\dfrac{z-z_1}{z-z_2}\right)>\dfrac\pi4$ where $|z|>1$}
\opt{B}{$|z_1z_2+z_2z_3+z_3z_1|=|z_1+z_2+z_3|$}
\opt{C}{$\mathrm{Im}\left(\dfrac{(z_1+z_2)(z_2+z_3)(z_3+z_1)}{z_1\cdot z_2\cdot z_3}\right)=0$}
\opt{D}{If $|z_1-z_2|=\sqrt2|z_1-z_3|=\sqrt2|z_2-z_3|$, then $\mathrm{Re}\left(\dfrac{z_3-z_1}{z_3-z_2}\right)=0$}
\end{qoptions}
\qn{8} From a point $(\sin\theta,\cos\theta)$ if three normals can be drawn to the parabola $y^2=4ax$ then the value of `$a$' belong to
\begin{qoptions}
\opt{A}{$\left(\dfrac12,1\right)$}
\opt{B}{$\left(-\dfrac12,0\right)$}
\opt{C}{$\left(0,\dfrac12\right)$}
\opt{D}{$(1,\infty)$}
\end{qoptions}
\qn{9} The number of ordered triplets $(x,y,z)$ of non negative integers satisfying the conditions. $x+y+z\le100$ and $x\le y\le z$
\begin{qoptions}
\opt{A}{if $x$ is odd is 14724}
\opt{B}{if $x$ is odd is 14722}
\opt{C}{if $x$ is even is 16065}
\opt{D}{if $x$ is even is 16164}
\end{qoptions}
\qn{10} All the five digit numbers in which each successive digit exceeds its predecessor are arranged in the increasing order. The 105th number does not contain the digit
\begin{qoptions}
\opt{A}{1}
\opt{B}{2}
\opt{C}{3}
\opt{D}{4}
\end{qoptions}
\qn{11} If $z_1=a+ib$ and $z_2=c+id$ are complex number such that $|z_1|=|z_2|=1$ and $\mathrm{Re}\left(z_1\overline{z_2}\right)=0$, then the pair of complex numbers $z_3=a+ic$ and $z_4=b+id$ satisfies :
\begin{qoptions}
\opt{A}{$|z_3|=1$}
\opt{B}{$|z_4|=1$}
\opt{C}{$\mathrm{Re}\left(z_3\overline{z_4}\right)=0$}
\opt{D}{$\mathrm{Im}\left(z_3\overline{z_4}\right)=0$}
\end{qoptions}
\qn{12} $y=mx$ bisects two distinct chords drawn from $(4,4)$ on $y^2=4x$ if
\begin{qoptions}
\opt{A}{$m=-1/2$}
\opt{B}{$m=0$}
\opt{C}{$m=1/2$}
\opt{D}{$m=1$}
\end{qoptions}
\sect{SECTION-II}
\qn{1} Find the minimum value of $\left(x_1-x_2\right)^2+\left(\sqrt{2-x_1^2}-\dfrac9{x_2}\right)^2$, where $x_1\in\left(0,\sqrt2\right)$ and $x_2\in R^+$.
\qn{2} Let 5 letters words are formed using the letters of the word `CALCULUS' and $N$ is the rank of word CALCU among them as arranged in dictionary. Then $N$ is
\qn{3} The radius of circle which passes through the focus of parabola $x^2=4y$ and touches it at point $(6,9)$ is $k\sqrt{10}$, then $k=$
\qn{4} A normal is drawn to the ellipse $\dfrac{x^2}{\left(a^2+2a+2\right)^2}+\dfrac{y^2}{\left(a^2+1\right)^2}=1$ whose centre is at origin. If maximum radius of the circle, centered at the origin and touching the normal is 5, then find the positive value of `$a$'.
\qn{5} If the normals at the four points $(x_1,y_1),(x_2,y_2),(x_3,y_3)$ and $(x_4,y_4)$ on the ellipse $\dfrac{x^2}{a^2}+\dfrac{y^2}{b^2}=1$ are concurrent, then find the value of $\left(x_1+x_2+x_3+x_4\right)\cdot\left(\dfrac1{x_1}+\dfrac1{x_2}+\dfrac1{x_3}+\dfrac1{x_4}\right)$.
\qn{6} The mean square deviation of a set of $n$ observations $x_1,x_2,\dots,x_n$ about a point $c$ is defined as $\dfrac1n\displaystyle\sum_{i=1}^n\left(x_i-c\right)^2$. The mean square deviation about $-2$ and 2 are 18 and 10 respectively, then standard deviation of this set of observations is :
\end{multicols}
\end{document}
